% This must be in the first 5 lines to tell arXiv to use pdfLaTeX, which is strongly recommended.
\pdfoutput=1
% In particular, the hyperref package requires pdfLaTeX in order to break URLs across lines.

\documentclass[11pt]{article}

% Change "review" to "final" to generate the final (sometimes called camera-ready) version.
% Change to "preprint" to generate a non-anonymous version with page numbers.

% CHANGER ENTRE MODE REVIEW ET PAS REVIEW !!

%\usepackage[review]{acl}
\usepackage{acl}
\usepackage{graphicx}


% Standard package includes
\usepackage{times}
\usepackage{latexsym}

% For proper rendering and hyphenation of words containing Latin characters (including in bib files)
\usepackage[T1]{fontenc}
% For Vietnamese characters
% \usepackage[T5]{fontenc}
% See https://www.latex-project.org/help/documentation/encguide.pdf for other character sets

% This assumes your files are encoded as UTF8
\usepackage[utf8]{inputenc}

% This is not strictly necessary, and may be commented out,
% but it will improve the layout of the manuscript,
% and will typically save some space.
\usepackage{microtype}

% This is also not strictly necessary, and may be commented out.
% However, it will improve the aesthetics of text in
% the typewriter font.
\usepackage{inconsolata}

% If the title and author information does not fit in the area allocated, uncomment the following
%
%\setlength\titlebox{<dim>}
%
% and set <dim> to something 5cm or larger.

\title{Using zero-shot prompting to generate synthetic DNN features for text classification.}

% Author information can be set in various styles:
% For several authors from the same institution:
% \author{Author 1 \and ... \and Author n \\
%         Address line \\ ... \\ Address line}
% if the names do not fit well on one line use
%         Author 1 \\ {\bf Author 2} \\ ... \\ {\bf Author n} \\
% For authors from different institutions:
% \author{Author 1 \\ Address line \\  ... \\ Address line
%         \And  ... \And
%         Author n \\ Address line \\ ... \\ Address line}
% To start a separate ``row'' of authors use \AND, as in
% \author{Author 1 \\ Address line \\  ... \\ Address line
%         \AND
%         Author 2 \\ Address line \\ ... \\ Address line \And
%         Author 3 \\ Address line \\ ... \\ Address line}

\author{Omar Cherif \\
  LORIA \\
  \texttt{mohamed.cherif@inria.fr} \\\And
  Christophe Cerisara \\
  LORIA\\
  \texttt{christophe.cerisara@loria.fr}  \\\AND
  Julien Falgas \\
  CREM\\
  Université de Lorraine \\
  \texttt{julien.falgas@univ-lorraine.fr} \\}

\begin{document}
\maketitle
\begin{abstract}
Large Language Models (LLMs) are capable of classifying documents from a simple prompt. In this article, we propose using an LLMs capacity to answer simple questions to obtain features to use in an MLP. We evaluate the created features on a classification task and by explore the implications of the results we reach the conclusion that logits encode valuable general information concerning the input. We obtain results equivalent to a Finetuned Q-LoRa model on a sincerity classification task.
\end{abstract}

\section{Introduction}

Feature Generation is a field of machine-learning that allows to get better results from classical algorithms, by tweaking and/or analyzing the input in order to get more valuable information out of it. \citep{florescuAlgorithmicallyGeneratingNew2019}. Large Language Models (LLMs) are able to answer simple questions concerning texts by giving a probability of the next token, these lists of probabilities are called logits. 

\section{Related works}


\section{Methodology}
\subsection{The classifications task}
To evaluate the quality of logits, we have mainly decided to use a subset of the Insincere Questions dataset\footnote{https://www.kaggle.com/competitions/quora-insincere-questions-classification/data}. To ensure a good distribution of the targets we have changed the dataset so that both classes are equally present (50/50) and then took a subset of 500 examples from that sample.

The same method has been applied to the Text document classification Dataset \footnote{https://www.kaggle.com/datasets/sunilthite/text-document-classification-dataset} and the IMDB movie review dataset \footnote{https://www.kaggle.com/datasets/lakshmi25npathi/imdb-dataset-of-50k-movie-reviews}.


\subsection{The questions}
To evaluate the quality of logits as features, we generate 5 sets of questions to which the answers will be the features:

\begin{itemize}
    \item Expert questions based on \cite{depauloCuesDeception2003}
    \item AI generated questions from Chat-GPT.
    \item Vocabulary questions concerning the text.
    \item Synesthesic questions concerning the text.
    \item Lists of random words.
\end{itemize}

\subsubsection{Expert questions:}
\cite{depauloCuesDeception2003} presents a list of cues extracted from a meta-study containing verbal and non-verbal cues. We have extracted only the verbal cues and transformed them into yes/no questions using an LLM. 
These questions are complex and require an in depth analysis and understanding of the text to answer them most of the time. 

\subsubsection{AI generated questions:}
We then generated a list of plausible questions to ask oneself if we try to determine if a question is sincere. We used Chat-GPT to generate a list of 20 yes/no questions. 

\subsubsection{Vocabulary questions:}
For this set of questions we generated 10 synonyms for the word sincere, and for each question took the form: "Is the previous text \begin{it} adjective \end{it}". And added the question asking whether it was sincere, totalling 11 features.

\subsubsection{Synesthetic questions:}
Synesthesia is a phenomenon where stimulation in one sensory modality automatically triggers a perception in a second modality \citep{cytowicSynesthesiaUnionSenses2002}. In this case since a large language model only perceives in one modality (text) there is no real synesthesia, but the questions can still be defined as synesthetic for humans. 
We have generated 20 questions that associated the text with human senses (smell, taste, touch, sight and hearing) using a Large Language Model. 

These control questions do not relate in any sense to the "sincerity" of a text.

\subsubsection{Gibberish texts}
Finally, gibberish texts consist of 10 lists of 8 random words generated using a website.

\subsection{Feature extraction}
By using a LLM (BloomZ-7B) \citep{muennighoffCrosslingualGeneralizationMultitask2023}, we present the text to be analyzed, followed by a question from one of the sets. Then, instead of taking the result, we look at the logits and extract the probability of "yes" minus "no". We end up with a dataframe containing probabilities for each questions.

\subsection{Classification}
We train the features on an MLP. We obtain the following results based on the features.

\section{Results}
We evaluate features individually and then look at how 
\begin{figure}
    \centering
    \includegraphics[width=1\linewidth]{images/VocabFeatures.png}
    \caption{The vocabulary features perform well together}
    \label{fig:enter-label}
\end{figure}
\section{Conclusion}

\begin{figure}
    \centering
    \includegraphics[width=1\linewidth]{images/Synestheticfeatures.png}
    \caption{The Synesthetic features perform better than expert questions.}
    \label{fig:enter-label}
\end{figure}
\subsection{In other tasks}
We also tried using the same synesthetic features for different tasks to see whether the information contained in the synesthetic answers could be used in a variety of situations. Notably, on the Text document classification dataset available on Kaggle which is based on the task of classifying small texts in 5 different genres (politics, sports, technology, entertainment and business) each new feature seems to help with the quality of the model as seen in Fig 3.
\begin{figure}
    \centering
    \includegraphics{images/classify.png}
    \caption{Caption}
    \label{fig:enter-label}
\end{figure}
Thus, it seems that a LLM is able to recognize if a text is about politics or technology by thinking about abstract features such as its "taste" its "smell" or its "feel".

\subsection{Discussion}
We argue that since LLMs perceive everything in a single modality, the notion of synesthesia (which can be conceived as "bridges" between modalities) is inaccurate. The model only knows about text and can thus answer questions about the "smell" of a text in an accurate way for classifications to be made.

\bibliography{Mabibliothèque,custom}

\appendix

\section{Questions}
\label{sec:appendix}

This is an appendix.

\end{document}
